\documentclass[book.tex]{subfiles}
\begin{document}

\chapter{Optimizing Mechanics Models in Cartesian Coordinates}
\label{chap:hamiltonian_nn}
\noindent\rule{11cm}{0.4pt} \\
\textbf{Prerequisites:} \autoref{chap:hamiltonians_mechanics}, \autoref{chap:neural_ode} \\
\textbf{Difficulty Level:} ** \\
\noindent\rule{11cm}{0.4pt}
\vspace{5mm}

The Hamiltonian and Lagrangian mechanics models we have learned about in previous chapters only hold for simple systems with an energy conservation law. However many real world mechanics problems require modeling of friction, contacts and collisions. Modeling such events introduces mathematical complexities since they introduce constraints and discontinuities into the system that standard differential equations cannot easily model. To solve such systems, we introduce a constrained optimizer into the differentiable program that models these systems  \cite{finzi2020generalizing}, \cite{zhong2021benchmarking}, \cite{zhong2021extending}.
%"""
%Abstract: The incorporation of appropriate inductive bias plays a critical role in learning dynamics from data. A growing body of work has been exploring ways to enforce energy conservation in the learned dynamics by encoding Lagrangian or Hamiltonian dynamics into the neural network architecture. These existing approaches are based on differential equations, which do not allow discontinuity in the states and thereby limit the class of systems one can learn. However, in reality, most physical systems, such as legged robots and robotic manipulators, involve contacts and collisions, which introduce discontinuities in the states. In this paper, we introduce a differentiable contact model, which can capture contact mechanics: frictionless/frictional, as well as elastic/inelastic. This model can also accommodate inequality constraints, such as limits on the joint angles. The proposed contact model extends the scope of Lagrangian and Hamiltonian neural networks by allowing simultaneous learning of contact and system properties. We demonstrate this framework on a series of challenging 2D and 3D physical systems with different coefficients of restitution and friction. The learned dynamics can be used as a differentiable physics simulator for downstream gradient-based optimization tasks, such as planning and control.
%"""


%"""
%The last few years have witnessed an increased interest in incorporating physics-informed inductive bias in deep learning frameworks. In particular, a growing volume of literature has been exploring ways to enforce energy conservation while using neural networks for learning dynamics from observed time-series data. In this work, we survey ten recently proposed energy-conserving neural network models, including HNN, LNN, DeLaN, SymODEN, CHNN, CLNN and their variants. We provide a compact derivation of the theory behind these models and explain their similarities and differences. Their performance are compared in 4 physical systems. We point out the possibility of leveraging some of these energy-conserving models to design energy-based controllers.
%"""

\section{Embedding the System into Cartesian Coordinates}

We first study methods to explicitly encode constraints into Hamiltonian and Lagrangian neural networks. A core technique is to embed a system description into Cartesian coordinates and then enforce a constraint via a Lagrangian multiplier. This technique simplifies the optimization difficulty for the problem.

%"""
%Reasoning about the physical world requires models that are endowed with the right inductive biases to learn the underlying dynamics. Recent works improve generalization for predicting trajectories by learning the Hamiltonian or Lagrangian of a system rather than the differential equations directly. While these methods encode the constraints of the systems using generalized coordinates, we show that embedding the system into Cartesian coordinates and enforcing the constraints explicitly with Lagrange multipliers dramatically simplifies the learning problem. We introduce a series of challenging chaotic and extended-body systems, including systems with N-pendulums, spring coupling, magnetic fields, rigid rotors, and gyroscopes, to push the limits of current approaches. Our experiments show that Cartesian coordinates with explicit constraints lead to a 100x improvement in accuracy and data efficiency.
%"""

We can model a system in Cartesian coordinates by a vector of Cartesian coordinates

\begin{align}
    x(t) &= (x_1(t), \dotsc, x_D(t))
\end{align}
where $D$ is the dimensionality of the system. The Lagrangian for this system is given by
\begin{align}
    L(x, \dot{x}) &= \frac{1}{2}\dot{x}^T M \dot{x}
\end{align}
where the constant mass matrix $M$ is given by
\begin{align}
    M&= \nabla_{\dot{x}}\nabla_{\dot{x}}^TL
\end{align}
Newton's second law for this system is given by %(\textcolor{red}{TODO: Add this form earlier})
\begin{align}
\delta x^T \left [ \frac{d}{dt}\frac{\partial L}{\partial \dot{x}} - \frac{\partial L}{\partial x} \right ] &= 0
\end{align}
%(\textcolor{red}{TODO: Provide intuition for the mass matrix}). 
If a system has $M$ degrees of freedom where $M < D$, then there must be $D - M$ constraints. We assume that all constraints are holonomic, %(\textcolor{red}{TODO: Explain holonomy/holonomic constraints in math section.}). 
which means that we can thus express the constraint in the form
\begin{align}
\Phi_k(x) &= 0
\end{align}
where $\Phi_k$ is some constraint function. The updated Newton's law is then given by
\begin{align}
\delta x^T \left [ \frac{d}{dt}\frac{\partial L}{\partial \dot{x}} - \frac{\partial L}{\partial x}   + (D_x \Phi)^T \lambda_L\right ] &= 0
\end{align}
where $\lambda_L$ is a Lagrange multiplier.

A Symplectic ODENet learns $M$, $V$
\begin{align}
\begin{pmatrix}
    \dot{q} \\
    \dot(p)
\end{pmatrix} &= \begin{pmatrix}
    \frac{\partial H}{\partial p} \\
    -\frac{\partial H}{\partial q}
\end{pmatrix}
\end{align}
\begin{align}
H(p, q) &= \frac{1}{2}M^{-1}(q)p + V(q)
\end{align}
\begin{align}
\dot{z} &= f_\theta(z) \\
\hat{z} &= \textrm{ODESolve}(z_{t_\theta}, f_\theta, t_1,\dotsc,t_n) \\
\mathcal{L} &= \|Z - \hat{Z} \|
\end{align}

\section{Limitations}

Collision and contact events cannot be easily modeled by Lagrangian and Hamiltonian networks.

\section{Constrained Lagrangian and Hamiltonian NNs}

%Generalized coordinates versus Cartesian coordinates. 
Cartesian coordinates are easier to learn than natural angular coordinates for a gyroscope or other systems of similar complexity. The Cartesian coordinates are referred to as constrained Hamiltonian and Lagrangian neural networks.


\end{document}
